\section{Accessibility by bounded subcomplexes}

We use CW complexes with a distinguished zero-cell as basepoint. The size of a CW complex in this section means its number of cells. This convention avoids confusing the cell cardinality of a CW complex with the generally larger cardinality of its singular simplicial set.

\begin{claim}{B.1}{Proved}{The bounded-subcomplex closure lemma}\label{B1}
Let $\mathcal A$ be a class of pointed CW homotopy types containing the point, closed under pointed homotopy colimits and all three two-out-of-three implications in homotopy cofiber sequences. Let $\kappa$ be an uncountable regular cardinal in $\mathcal U$. Define $\mathcal A^\kappa$ by the condition that every pointed subcomplex of fewer than $\kappa$ cells is contained in a pointed subcomplex of fewer than $\kappa$ cells belonging to $\mathcal A$.

This defines a class of homotopy types, is contained in $\mathcal A$, contains every member of $\mathcal A$ with fewer than $\kappa$ cells, and has the same closure properties. It is closed under desuspension and retracts. It is generated under pointed homotopy colimits by a $\mathcal U$-small set, namely representatives of the members of $\mathcal A$ with fewer than $\kappa$ cells.
\end{claim}
\begin{proof}
We first record the cardinal bookkeeping. The image of a subcomplex with fewer than $\kappa$ cells under a cellular map is contained in a subcomplex with fewer than $\kappa$ cells. For each closed cell its image is compact and hence lies in a finite subcomplex; taking their union and all faces gives the assertion. The same statement holds for a cellular homotopy, since each closed cell times the unit interval is compact. Countably many repetitions still involve fewer than $\kappa$ cells. Regularity also permits taking unions of fewer than $\kappa$ such subcomplexes.

For $X\in\mathcal A^\kappa$, let $P_X$ be the poset of its small pointed $\mathcal A$-subcomplexes. This poset is $\kappa$-filtered: the union of fewer than $\kappa$ objects is small and can be enlarged by the defining property. Its union is $X$. This is also its homotopy colimit. Indeed a sphere or homotopy in the union meets finitely many cells and hence lands in a common member. In the bar model for the homotopy colimit it meets finitely many cells and index strings; a common upper bound contracts this finite support into a single member. These two observations prove that the natural comparison induces isomorphisms on all homotopy groups. Thus $X\in\mathcal A$. A small $X\in\mathcal A$ satisfies the condition by taking $X$ itself.

We next prove independence of the CW model. Let $f:X\to Y$ and $g:Y\to X$ be pointed cellular homotopy inverses, and choose cellular homotopies from $gf$ and $fg$ to the respective identities. Suppose $X\in\mathcal A^\kappa$ and start with a prescribed small pointed subcomplex $Y_0\subseteq Y$. Repeatedly enlarge small subcomplexes on the two sides so that the next ones contain the images under $f,g$ and the tracks of both homotopies of the preceding ones. At each stage additionally enlarge the $X$-side to a small member of $\mathcal A$. The cardinal bookkeeping above permits all choices. At the countable union we obtain small subcomplexes $X_\infty,Y_\infty$, containing the initial data, on which $f,g$ and their inverse homotopies restrict. The chosen $\mathcal A$-subcomplexes on the $X$-side are cofinal in its sequential union, so $X_\infty\in\mathcal A$. Hence $Y_\infty\in\mathcal A$ as well. This proves $Y\in\mathcal A^\kappa$. A pointed weak equivalence of CW complexes is a pointed homotopy equivalence, so the class is well defined on homotopy types.

For closure under pointed homotopy colimits, use the cellular bar model of a strict pointed diagram $F:I\to\mathrm{CW}_*$. Its cells are indexed by finite strings of arrows of $I$ and cells of the corresponding $F(i)$; faces use the structure maps. A functorial realization of singular complexes, if needed, makes all structure maps cellular without changing their homotopy types, and the preceding paragraph transfers $\mathcal A^\kappa$ to those models. Begin with a small subcomplex of the bar realization. The cells it meets specify fewer than $\kappa$ objects, arrows, and cells in the diagram. The subcategory $J\subseteq I$ generated by those arrows has fewer than $\kappa$ arrows, since its morphisms are finite words in the chosen arrows.

For every object of $J$, enlarge the prescribed small subcomplex of $F(i)$ to a small $\mathcal A$-subcomplex. Enlarge again to include all images under arrows in $J$, then again to small $\mathcal A$-subcomplexes, and repeat. There are fewer than $\kappa$ objects and arrows, and only countably many rounds. The resulting unions $F'(i)\subseteq F(i)$ are small members of $\mathcal A$ and form a $J$-subdiagram. Their pointed bar construction is a CW subcomplex of the original bar construction: each face remains in the chosen diagram, degeneracies add no new cells, and the basepoint section is a subcomplex before it is collapsed. It contains the prescribed cells, has fewer than $\kappa$ cells, and belongs to $\mathcal A$ by homotopy-colimit closure. This proves the required local smallness of the original homotopy colimit.

We give the simultaneous cofiber argument in detail. Model a pointed cofiber sequence by a CW subcomplex inclusion $X\subseteq Y$ and quotient $C=Y/X$, with quotient map $q:Y\to C$. For a small subcomplex $Y_0\subseteq Y$, put $X_0=Y_0\cap X$ and $C_0=q(Y_0)$; then $C_0=Y_0/X_0$ with its inherited CW structure. A prescribed small subcomplex of $X$ can be included by adding its cells to $Y_0$. A prescribed small subcomplex of $C$ can be included by lifting its non-basepoint cells to their unique cells of $Y\setminus X$ and adding their boundaries in $Y$. These boundaries meet only finitely many cells for each characteristic map. This yields a small larger $Y_1$, still with $X_1=Y_1\cap X$ and $C_1=Y_1/X_1$.

Suppose any chosen two of $X,Y,C$ belong to $\mathcal A^\kappa$. At stage $r$, enlarge the small subcomplex in each of those two vertices to a small member of $\mathcal A$, and then restore the three compatibility equalities by the preceding operations. Repeat countably. At the union the triple is still
\[
 X_\infty=Y_\infty\cap X,
 \qquad C_\infty=q(Y_\infty)=Y_\infty/X_\infty.
\]
In each of the two designated vertices, the chosen $\mathcal A$-subcomplexes are cofinal in that union. Both vertices therefore belong to $\mathcal A$. Its cofiber two-out-of-three property gives the third. Starting with any prescribed small subcomplex of the third vertex, lifted if it is in $C$, proves the defining property for that vertex. This treats all three implications, in particular the sometimes omitted implication from $Y,C$ to $X$.

Here is also a direct desuspension check. If $\Sigma X\in\mathcal A^\kappa$, begin with a small $X_0\subseteq X$ and choose successively
\[
 \Sigma X_0\subseteq K_0\subseteq\Sigma X_1
       \subseteq K_1\subseteq\Sigma X_2\subseteq\cdots,
\]
where each $K_r$ is a small member of $\mathcal A$. A small subcomplex of a cellular suspension is contained in the suspension of a small subcomplex, so the intervening $X_r$ exist. At the union, $\Sigma X_\infty=\bigcup_rK_r\in\mathcal A$; cofiber two out of three for $X_\infty\to *\to\Sigma X_\infty$ gives $X_\infty\in\mathcal A$. Its size is still less than $\kappa$. Thus $X\in\mathcal A^\kappa$.

Retracts follow from homotopy-colimit closure: if $Y$ is a pointed homotopy retract of $X$, the mapping telescope of the corresponding homotopy idempotent of $X$ is equivalent to $Y$. Finally all CW complexes with fewer than $\kappa$ cells have underlying sets of cardinality at most $\max(\kappa,2^{\aleph_0})$. There is a $\mathcal U$-small set of their CW structures and maps up to isomorphism, hence also a set of their homotopy-type representatives. The filtered presentation by $P_X$ and the already proved closure show that the representatives lying in $\mathcal A$ generate exactly $\mathcal A^\kappa$ under pointed homotopy colimits.
\end{proof}

\begin{claim}{B.2}{Proved}{An accessible reversible sublocalizer}\label{B2}
With the hypotheses of B.1, the class
\[
 W^\kappa_{\mathcal A}
   =\{u:C(Bu)\in\mathcal A^\kappa\}
\]
is an accessible reversible fundamental localizer. A generating set consists of category-of-simplices models of all maps between CW complexes with fewer than $\kappa$ cells whose homotopy cofibers belong to $\mathcal A$.
\end{claim}
\begin{proof}
B.1 and A.3 prove that it is a reversible fundamental localizer. There is a $\mathcal U$-small set $S_\kappa$ of representatives of the indicated arrows, by the last paragraph of B.1. Applying the singular functor and then its category of simplices gives actual functors in $\Cat$; the natural realization comparisons preserve the specified arrows and their cofibers. The resulting set is still $\mathcal U$-small, even though its categories may have more cells or simplices than the original CW complexes. Every member is in $W^\kappa_{\mathcal A}$.

Conversely let a map $f$ have $C(f)\in\mathcal A^\kappa$. If its source is empty and its target is nonempty, $C(f)=Y_+$ has $S^0$ as a pointed retract. Then $S^0\in\mathcal A^\kappa$, and pointed homotopy-colimit closure makes this class all pointed spaces: cells build every pointed CW complex from $S^0$. The localizer is total, generated by $\varnothing\to *$. The empty identity causes no difficulty. We may therefore assume the source is nonempty. Replace $f$ by the inclusion $X\subseteq M_f$ of one end of its cellular mapping cylinder; its target deformation retracts onto the original target. Put $C=M_f/X$ and include a fixed source point in every subpair. Given any small subpair, its quotient image is a small pointed subcomplex of $C$. Enlarge this image to a small $C'\in\mathcal A$ and lift the noncollapsed cells of $C'$ to $M_f$, including their boundaries and the prescribed source cells. If $Y'$ is the resulting subcomplex and $X'=Y'\cap X$, then
\[
 Y'/X'=C',
\]
not merely a space mapping to $C'$. Thus $X'\hookrightarrow Y'$ is a small arrow belonging to $S_\kappa$ up to its chosen model. Such subarrows are $\kappa$-filtered and cofinal among all small subarrows of $X\subseteq M_f$. Their colimit, and their homotopy colimit as an arrow, is $X\subseteq M_f$.

Every fundamental localizer is closed under homotopy colimits of arrows of spaces after passage through categories of simplices. This is the relative-Theorem-A argument checked in A.3, or explicitly the Grothendieck-construction argument of POLES B.2. Hence every fundamental localizer containing $S_\kappa$ contains $f$. It follows that $W^\kappa_{\mathcal A}$ is exactly the fundamental localizer generated by the stated small set.

The use of arbitrary \emph{arrows} in this generating set matters. Replacing them by the object-to-point maps would generally produce the proper nullification with the same aspherical spaces, and would not prove the asserted equality of localizers.
\end{proof}

\begin{claim}{B.3}{Proved}{The accessible question has a negative stable-classification answer}\label{B3}
In ZFC, there exists an accessible reversible fundamental localizer that is not $R(\mathcal L)$ for any localizing subcategory $\mathcal L$ of spectra. One example is $W^{\aleph_1}_{\Nul}$.
\end{claim}
\begin{proof}
The countable CW complex $\RP^\infty$ belongs to $\Nul$ by A.4, hence to $\Nul^{\aleph_1}$ by B.1. The space $\RP^2$ does not belong to $\Nul$, and so does not belong to its subclass $\Nul^{\aleph_1}$. B.2 makes the associated localizer accessible and reversible. A.5 shows that any localizer of the form $R(\mathcal L)$ containing the first object-to-point map must contain the second. This contradiction proves the assertion. The construction uses a fixed regular cardinal and explicit countable enlargement arguments; no form of Vop\v enka's principle is used.
\end{proof}

\begin{claim}{B.4}{Proved}{Reversible closure of a small set is accessible}\label{B4}
The smallest reversible fundamental localizer containing any $\mathcal U$-small set of functors is accessible. In particular the smallest reversible localizer containing $\RP^\infty\to *$ is accessible and is not $RT$-fixed.
\end{claim}
\begin{proof}
Take the intersection $W$ of the reversible fundamental localizers containing the chosen set $S$, and let $\mathcal A$ be its aspherical pointed spaces. Intersections preserve all the required axioms. The cofiber criterion and forward homotopy-colimit closure imply that $\mathcal A$ has the hypotheses of B.1: pointed homotopy colimits of aspherical spaces are aspherical, and the cofiber criterion supplies all three cofiber implications. Choose an uncountable regular $\kappa\in\mathcal U$ exceeding the cell cardinalities of models for the cofibers of every member of $S$. Then $W^\kappa_{\mathcal A}\subseteq W$ and it contains $S$. It is reversible by B.2, so minimality gives the reverse inclusion. Thus $W=W^\kappa_{\mathcal A}$ is accessible.

For $S=\{\RP^\infty\to *\}$, A.3--A.4 give $W\subseteq W_S$ and therefore exclude $\RP^2\to *$. A.5 puts that map in $RT(W)$. All intersections may be taken as sets in the next universe; the resulting generating set in B.2 belongs to $\mathcal U$. Equivalently one can define the smallest closure by well-founded set-sized derivation trees for the localizer and reverse-square rules, which gives the same separation-based construction in ordinary class notation.
\end{proof}
